% GENERATED FILE. Edit canonical lessons, then run npm run generate:books.
% canonical-source-hash: fnv1a64:7920f9806dcd2242
% canonical-lessons: ZH-C539-re-provoke, ZH-C539-guai-blame, ZH-W539-guai, ZH-C539-guaishi, ZH-C539-chikui, ZH-C539-chuanghuo

\chapter{Whose Fault?}
\label{ch:zh-whose-fault}
\hlchaptermodality{539}

\begin{chapteropening}
\textbf{By the end of this chapter:} \emph{I can say who is to blame, say something is strange, and talk about provoking someone, coming off worse and getting into trouble.}

Complete the last lesson of chapter 539: I can say who is to blame, say something is strange, and talk about provoking someone, coming off worse and getting into trouble.
\end{chapteropening}

\section[rě]{rě — to stir up trouble; to provoke}

\label{lesson:ZH-C539-re-provoke}



\begin{quote}
Before the new one: say the Mandarin for to arise from something, then the Mandarin for not in the least.
\end{quote}

\subsection*{You'll want to know: rě}
\emph{rě} — \textquotedblleft{}to stir up trouble; to provoke\textquotedblright{}.

\begin{quote}
\emph{nǐ búyào rě tā bùrán tā huì shēngqì} — Don't provoke him, or he'll get angry.
\end{quote}

Its characters come later; for now it is a word for the ear.

\begin{cousinweb}
\emph{rě} is inherited from Middle Chinese \emph{nyaeX} (Guangyun, in Baxter's transcription), glossed 'a troubled mind'; the dictionary also lists \emph{nyak}. The sense 'to provoke' is not given there.
\end{cousinweb}

\begin{grammarlens}[title={What you've built}]
One more word for saying what would happen.
\end{grammarlens}

\subsection*{Your turn}
\begin{itemize}
\raggedright
  \item \emph{Say it:} \emph{rě}
  \item \emph{Say it:} \emph{rě}, once more
  \item \emph{Say it:} \emph{nǐ búyào rě tā bùrán tā huì shēngqì} — Don't provoke him, or he'll get angry.
  \item \emph{Recall:} say the Mandarin for to arise from something, then the Mandarin for not in the least, then say \emph{rě} again
\end{itemize}

\begin{tcolorbox}[breakable,colback=teal!4,colframe=teal!35!black,title={Before you move on}]
What does rě mean? (\textquotedblleft{}To stir up trouble; to provoke\textquotedblright{}.) Say it once more.
\end{tcolorbox}
\section[\texorpdfstring{\zh{怪} (guài)}{怪 (guài)}]{\zh{怪} (guài) — to put the blame on someone}

\label{lesson:ZH-C539-guai-blame}



\begin{quote}
Before the new one: say the Mandarin for not in the least, then the Mandarin for to stir up trouble.
\end{quote}

\subsection*{You'll want to know: \zh{怪}}
\textbf{\zh{怪}} — \emph{guài} — \textquotedblleft{}to put the blame on someone\textquotedblright{}.

\begin{quote}
\textbf{\zh{这不能怪你原因不在你}} \emph{zhè bù néng guài nǐ yuányīn bú zài nǐ} — You can't be blamed for this: the cause doesn't lie with you.
\end{quote}

The new character, in pieces: \zh{怪} — eight strokes in three pieces: the upright heart on the left, three strokes, as in \zh{快}; on the right \zh{又}, two strokes; and under it the earth piece, three strokes. The next lesson puts it in your hand.

\begin{cousinweb}
\textbf{\zh{怪}} (\emph{guài}) is inherited from Middle Chinese \emph{kweajH} (Guangyun, in Baxter's transcription), glossed 'strange'. The sense 'to blame' is not given there; it is usually explained as finding fault, but how it came is not settled.
\end{cousinweb}

\begin{grammarlens}[title={What you've built}]
One more word for saying what would happen.
\end{grammarlens}

\subsection*{Your turn}
\begin{itemize}
\raggedright
  \item \emph{Say it:} \emph{guài}
  \item \emph{Say it:} \emph{guài}, once more
  \item \emph{Say it:} \emph{zhè bù néng guài nǐ yuányīn bú zài nǐ} — You can't be blamed for this: the cause doesn't lie with you.
  \item \emph{Recall:} say the Mandarin for not in the least, then the Mandarin for to stir up trouble, then say \emph{guài} again
\end{itemize}

\begin{tcolorbox}[breakable,colback=teal!4,colframe=teal!35!black,title={Before you move on}]
What does \zh{怪} mean? (\textquotedblleft{}To put the blame on someone\textquotedblright{}.) Say it once more.
\end{tcolorbox}
\section[\texorpdfstring{\zh{怪} (guài)}{怪 (guài)}]{\zh{怪} — the upright heart, then \zh{又} over the earth piece}

\label{lesson:ZH-W539-guai}



\begin{quote}
The word first, then the character inside it.

\begin{itemize}
\raggedright
  \item \emph{Recall:} say \emph{guài} — \textbf{R1}, one lesson back, before it decays
\end{itemize}

\begin{quote}
\zh{快} \zh{地}
\end{quote}

\emph{Read it:} each one aloud, with nothing said first — \emph{kuài}, \emph{dì}
\end{quote}

\subsection*{Script — the upright heart, then \zh{又} over the earth piece}
\hlblockfigure{\detokenize{figures/ZH-W539-guai-filmstrip.pdf}}{How \zh{怪} is written, stroke by stroke}

\begin{quote}
\zh{怪}
\end{quote}

\textbf{Three pieces: the upright heart on the left, then \zh{又}, then the earth piece under it.} Count the pieces, not the strokes.

\textbf{The upright heart, three strokes}, as in \zh{快}: a short dot on the left, a short dot on the right, then a long vertical down between them.

\textbf{\zh{又}, two strokes}, as you wrote it: a bar that turns into a long stroke falling to the left, then a long stroke falling to the right across it.

\textbf{The earth piece, three strokes}, as on the left of \zh{地}, but at the foot: a short bar, a vertical down through it, then a long bar along the foot.

3 + 2 + 3 = eight.

\subsection*{Your turn}
\begin{itemize}
\raggedright
  \item \emph{Write it:} \zh{怪}, all eight, with the model in front of you
  \item \emph{Cover:} write it again from memory, and count the pieces, not the strokes
\end{itemize}

\begin{tcolorbox}[breakable,colback=teal!4,colframe=teal!35!black,title={Before you move on}]
How many pieces does \zh{怪} have? (\textbf{Three: the upright heart, \zh{又} and the earth piece}.) Which piece is written last? (\textbf{The earth piece, at the foot}.)

Source: \href{https://raw.githubusercontent.com/chanind/hanzi-writer-data/68d10a4b21150cae5e1ebbd223eed289cf32d90c/data/\%E6\%80\%AA.json}{Hanzi Writer Data \zh{怪}, pinned stroke paths and medians}.
\end{tcolorbox}
\section[\texorpdfstring{\zh{怪事} (guàishì)}{怪事 (guàishì)}]{\zh{怪事} (guàishì) — a strange business; an odd thing}

\label{lesson:ZH-C539-guaishi}



\begin{quote}
Before the new one: say the Mandarin for to stir up trouble, then the Mandarin for to put the blame on.
\end{quote}

\subsection*{You'll want to know: \zh{怪事}}
\textbf{\zh{怪事}} — \emph{guàishì} — \textquotedblleft{}a strange business; an odd thing\textquotedblright{}.

\begin{quote}
\textbf{\zh{这真是怪事}} \emph{zhè zhēn shì guàishì} — That really is a strange business.
\end{quote}

\begin{cousinweb}
\textbf{\zh{怪事}} (\emph{guàishì}) is built from \zh{怪} 'strange' and \zh{事} 'a matter', both inherited from Middle Chinese, \emph{kweajH} and \emph{dzriH} (Guangyun, in Baxter's transcription), which glosses the first 'strange'.
\end{cousinweb}

\begin{grammarlens}[title={What you've built}]
One more word for saying what would happen.
\end{grammarlens}

\subsection*{Your turn}
\begin{itemize}
\raggedright
  \item \emph{Say it:} \emph{guàishì}
  \item \emph{Say it:} \emph{guàishì}, once more
  \item \emph{Say it:} \emph{zhè zhēn shì guàishì} — That really is a strange business.
  \item \emph{Recall:} say the Mandarin for to stir up trouble, then the Mandarin for to put the blame on, then say \emph{guàishì} again
\end{itemize}

\begin{tcolorbox}[breakable,colback=teal!4,colframe=teal!35!black,title={Before you move on}]
What does \zh{怪事} mean? (\textquotedblleft{}A strange business; an odd thing\textquotedblright{}.) Say it once more.
\end{tcolorbox}
\section[chīkuī]{chīkuī — to come off worse; to lose out}

\label{lesson:ZH-C539-chikui}



\begin{quote}
Before the new one: say the Mandarin for to put the blame on, then the Mandarin for a strange business.
\end{quote}

\subsection*{You'll want to know: chīkuī}
\emph{chīkuī} — \textquotedblleft{}to come off worse; to lose out\textquotedblright{}.

\begin{quote}
\emph{rúguǒ nǐ bù tīng wǒ de jiù huì chīkuī} — If you don't listen to me, you'll come off worse.
\end{quote}

Its characters come later; for now it is a word for the ear.

\begin{cousinweb}
\emph{chīkuī} is built from \zh{吃} 'to eat' and \emph{kuī} 'a loss'. For the first the dictionary has \emph{khek}, glossed 'to eat', under another spelling; the second is inherited from Middle Chinese \emph{khjwe} (Guangyun, in Baxter's transcription), glossed 'to be lacking': to swallow a loss.
\end{cousinweb}

\begin{grammarlens}[title={What you've built}]
One more word for saying what would happen.
\end{grammarlens}

\subsection*{Your turn}
\begin{itemize}
\raggedright
  \item \emph{Say it:} \emph{chīkuī}
  \item \emph{Say it:} \emph{chīkuī}, once more
  \item \emph{Say it:} \emph{rúguǒ nǐ bù tīng wǒ de jiù huì chīkuī} — If you don't listen to me, you'll come off worse.
  \item \emph{Recall:} say the Mandarin for to put the blame on, then the Mandarin for a strange business, then say \emph{chīkuī} again
\end{itemize}

\begin{tcolorbox}[breakable,colback=teal!4,colframe=teal!35!black,title={Before you move on}]
What does chīkuī mean? (\textquotedblleft{}To come off worse; to lose out\textquotedblright{}.) Say it once more.
\end{tcolorbox}
\section[chuǎnghuò]{chuǎnghuò — to get into trouble; to cause a mess}

\label{lesson:ZH-C539-chuanghuo}



\begin{quote}
Before the new one: say the Mandarin for a strange business, then the Mandarin for to come off worse.
\end{quote}

\subsection*{You'll want to know: chuǎnghuò}
\emph{chuǎnghuò} — \textquotedblleft{}to get into trouble; to cause a mess\textquotedblright{}.

\begin{quote}
\emph{yàoshi nǐ zài chuǎnghuò wǒ jiù gàosu nǐ māma} — If you get into trouble again, I'll tell your mum.
\end{quote}

Its characters come later; for now it is a word for the ear.

\begin{cousinweb}
\emph{chuǎnghuò} is built from \emph{chuǎng} 'to rush' and \emph{huò} 'disaster'. The second is inherited from Middle Chinese \emph{hwaX} (Guangyun, in Baxter's transcription), glossed 'harm'. The dictionary lists the first as \emph{trhimH}, a horse dashing out of a gate, but chuǎng is not the regular outcome of that form, so its history is not settled.
\end{cousinweb}

\begin{grammarlens}[title={What you've built}]
One more word for saying what would happen.
\end{grammarlens}

\subsection*{Your turn}
\begin{itemize}
\raggedright
  \item \emph{Say it:} \emph{chuǎnghuò}
  \item \emph{Say it:} \emph{chuǎnghuò}, once more
  \item \emph{Say it:} \emph{yàoshi nǐ zài chuǎnghuò wǒ jiù gàosu nǐ māma} — If you get into trouble again, I'll tell your mum.
  \item \emph{Recall:} say the Mandarin for a strange business, then the Mandarin for to come off worse, then say \emph{chuǎnghuò} again
\end{itemize}

\begin{tcolorbox}[breakable,colback=teal!4,colframe=teal!35!black,title={Before you move on}]
What does chuǎnghuò mean? (\textquotedblleft{}To get into trouble; to cause a mess\textquotedblright{}.) Say it once more.
\end{tcolorbox}
